% 一次掷骰的六个结果分别经点数X与奇偶指示变量I映射为数值。公平骰子，每个结果概率1/6。
\begin{tikzpicture}[line width=.6pt,font=\fontsize{8}{10}\selectfont,>=Stealth]
\tikzset{res/.style={circle,draw=FigureStroke,fill=white,minimum size=5.2mm,inner sep=0pt},
  val/.style={rounded corners=1.8pt,draw=FigureBlue,fill=FigureBlueFill,minimum width=6mm,minimum height=5mm,inner sep=1pt},
  ival/.style={rounded corners=1.8pt,draw=FigureGreen,fill=FigureGreenFill,minimum width=6mm,minimum height=5mm,inner sep=1pt},
  mapx/.style={->,draw=FigureBlue,line width=.6pt},
  mapi/.style={->,draw=FigureGreen,line width=.6pt}}
% 样本空间
\draw[FigureStroke,rounded corners=2pt] (-0.55,-0.45) rectangle (0.55,5.45);
\node[FigureInk] at (0,5.75) {样本空间 $\Omega$};
\foreach \k in {1,...,6}{\node[res] (w\k) at (0,{6-\k-0.0}) {$\k$};}
% X 的取值
\node[FigureInk] at (-3.0,5.75) {点数 $X(\omega)=\omega$};
\foreach \k in {1,...,6}{\node[val] (x\k) at (-3.0,{6-\k}) {$\k$}; \draw[mapx] (w\k) -- (x\k);}
% I 的取值
\node[FigureInk] at (3.2,5.75) {奇偶指示 $I(\omega)$};
\node[ival] (i1) at (3.2,3.75) {$1$};
\node[ival] (i0) at (3.2,1.25) {$0$};
\foreach \k in {2,4,6}{\draw[mapi] (w\k) -- (i1);}
\foreach \k in {1,3,5}{\draw[mapi] (w\k) -- (i0);}
\node[FigureMuted,font=\fontsize{7.5}{9}\selectfont,align=left,anchor=west] at (3.65,3.75) {$\mathbb{P}(I=1)$\\$=\mathbb{P}(\{2,4,6\})$};
\node[FigureMuted,font=\fontsize{7.5}{9}\selectfont,align=left,anchor=west] at (3.65,1.25) {$\mathbb{P}(I=0)$\\$=\mathbb{P}(\{1,3,5\})$};
\end{tikzpicture}
